%\pagenumbering{arabic}
\section{Preliminaries}

In this section, we introduce the notation and terminology used throughout the paper. We first recall diagram categories indexed by finite posets and the notion of universal derived equivalence. We then recall Ladkani's construction, which provides the main background for the constructions developed in the following section.

For general background on derived categories and homological algebra, see for instance  Weibel \cite{weibel} and Miličić \cite{dercat}.
For module categories, derived equivalences are classically studied through tilting theory and Rickard's derived Morita theory \cite{Rickard}.

\subsection{Diagram categories and universal derived equivalence}

Let $X$ be a finite poset. We regard $X$ as a category whose objects are the elements of $X$, and in which there is a unique morphism $x \to x'$ whenever $x \leq x'$.

The height $h(X)$ of a finite poset $X$ is the largest integer $r$ such that $X$ contains a chain $x_0<x_1< \dots <x_r$ of $r+1$ distinct elements.

If $\mathcal{A}$ is an abelian category, we denote by $\mathcal{A}^X$ the category of $X$-shaped diagrams in $\mathcal{A}$, that is, the functor category $\operatorname{Fun}(X,\mathcal{A})$. Thus, an object of $\mathcal{A}^X$ consists of a family $(A_x)_{x \in X}$ of objects of $\mathcal{A}$, together with morphisms $A_x \to A_{x'}$ for every relation $x \le x'$, compatible  with composition. Since kernels and cokernels in $\mathcal{A}^X$ are computed pointwise, $\mathcal{A}^X$ is abelian whenever $\mathcal{A}$ is abelian. We write $D(\mathcal{A}^X)$ for its derived category. 

Two finite posets $X$ and $Y$ are said to be universally derived equivalent \cite{ladkani} if, for every abelian category $\mathcal{A}$, there is an equivalence of triangulated categories
\[D(\mathcal{A}^X) \simeq D(\mathcal{A}^Y).\] 

Universal derived equivalence is therefore a combinatorial property of the indexing posets, independent of the choice of coefficients or of a particular category of representations.
\subsection{Ladkani's construction} \label{Ladkani's construction}

Given two posets $X$ and $Y$, for $y \in Y$ we define \[[y, \bullet ] \coloneq \{ y' \in Y \,\,\text{such that}\,\,y \leq y' \},\]
\[[\bullet, y ] \coloneq \{ y' \in Y \,\,\text{such that}\,\,y'\leq y \}.\]

Let $X$ and $Y$ be two finite posets and $\{Y_x\}_{x \in X}$ a collection of subsets of $Y$ such that the following properties hold.

\begin{enumerate}[label=(\roman*), ref=(\roman*)]
\item \label{prop:disjoint} $\forall x \in X, \,\forall y \neq y' \in Y_x $, \[
     [y, \bullet ] \cap [y', \bullet ]= \emptyset \text{\,\,\,\,and\,\,\,\,} [\bullet, y ] \cap [\bullet, y' ]= \emptyset;
\]
   
    
    \item \label{prop:phi} $\forall x,x' \in X$ such that $x \leq x'$, we have a bijection
    $\varphi_{xx'} \colon Y_x \to Y_{x'}$ such that $\forall y \in Y_{x}$, 
    $
    y \leq \varphi_{xx'}(y)$.
\end{enumerate}



Now, we define two relations $\leq_+$ and $\leq_-$ on the disjoint union $X \sqcup Y$.
These relations coincide with the original orders on $X$ and on $Y$ (i.e. if $x\leq x'$ in $[X,\leq]$, then $x \leq_+x'$ and $x \leq_- x'$, and similarly for elements inside $[Y,\leq]$).
Moreover, if $x\in X$ and $y \in Y$, we define:
\[x \leq_+y \iff \exists \,  y_x \in Y_x \text{\,\,such that\,\,}y_x \leq y,\]
\[y \leq_-x \iff \exists \,  y_x \in Y_x \text{\,\,such that\,\,}y \leq y_x.\]
These are the only mixed relations; in particular, for $x \in X$ and $y \in Y$, we never have $y \leq_+x$ or $x \leq_-y$.

Let $\mathcal{Y}=(Y_x)_{x \in X}$. We denote $X \sqcup Y$, equipped with the relations $\leq_+$ and $\leq_-$, by 
\[(X \sqcup Y)^+_{\mathcal{Y}}\,\,\,\,\, \text{and}\,\,\,\,\,(X \sqcup Y)^-_{\mathcal{Y}},\] respectively. When the family $\mathcal{Y}$ is clear from the context, we omit the subscript.

\begin{lemma}\label{lemma part ord}\cite[Section 1]{ladkani}
    The relations $\leq_+$ and $\leq_-$ defined above are partial orders on $X \sqcup Y$. Hence  $(X \sqcup Y)^+_{\mathcal{Y}}$ and $(X \sqcup Y)^-_{\mathcal{Y}}$ are posets. 
\end{lemma}
We call $(X \sqcup Y)_{\mathcal{Y}}^+$ and $(X \sqcup Y)^-_{\mathcal{Y}}$ the positive and negative posets associated with $\mathcal{Y}$, respectively.

\begin{theorem}\label{ladkani}

\cite[Theorem~1.1]{ladkani}
The positive and negative posets $(X \sqcup Y)^+_{\mathcal{Y}}$ and $(X \sqcup Y)^-_{\mathcal{Y}}$ are universally derived equivalent.

\end{theorem}

If each $Y_x$ is a singleton, say $Y_x=\{f(x)\}$, then the family $\{Y_x\}_{x \in X}$ is uniquely determined by a function $f \colon X \to Y$. In this case, assumption \cref{prop:disjoint} is automatic, while assumption \cref{prop:phi} is equivalent to $f$ being order-preserving. Thus, Theorem \ref{ladkani} specializes to the following result.

\begin{proposition}\cite[Corollary~1.3]{ladkani}\label{Ladkani singleton}
    Let $X$ and $Y$ be finite posets, and let $f \colon X \to Y$ be an order-preserving map. Define two partial orders $\leq_f^+$ and $\leq_f^-$ on the disjoint union $X \sqcup Y$ by requiring that they restrict to the given orders on $X$ and $Y$, and that, for $x \in X$ and $y \in Y$, 
    \[x \leq_f^+ y \iff f(x) \leq y, \,\,\,\,\,\,\,\, y \leq_f^- x \iff y \leq f(x).\]
    Then the two posets $(X \sqcup Y, \leq_f^+)$ and $(X \sqcup Y, \leq_f^-)$ are universally derived equivalent.
\end{proposition}

Ladkani later referred to this order-preserving map construction as a flip-flop \cite{ladkani2}; see also \cite[Section 5]{Goguet} for a recent use of this terminology.